%% ===== Discussion =====
\section{Discussion}
\label{sec:discussion}

% [29/10] So với bản 28/10: nêu tiêu chuẩn rút từ định nghĩa migration để rõ "gap" là so với cái gì; bỏ các câu rào chung ("observations about the selected corpus…") vì mỗi gap đã tự nêu giới hạn của dữ liệu. Sáu gap thành năm (xem lời dẫn 7.1). Pilot thêm bằng chứng cho G1–G3 bằng một tổ hợp oracle không bài nào trong map dùng (Bảng 5: O-h không đi cùng O-c hay O-f; đã kiểm hai bài O-h với bài gốc).
\new{RQ5 asks what gaps separate current evaluation practice from the needs of real migration projects, and what studies should report and investigate next. A migration replaces an existing program, so it needs evidence that the migrated code behaves like the original. Against this need, the map shows five gaps (Section~\ref{sec:rq5}). For three of them, the pilot adds evidence from a combination of oracles that no paper in the map uses: compilation, similarity, an LLM judge and existing tests, all applied to the same translations. The gaps lead to directions for further study, implications for three audiences (Section~\ref{sec:implications}) and a reporting checklist (Section~\ref{sec:checklist}).}

\subsection{RQ5: Gaps between evaluation practice and real migration}
\label{sec:rq5}

% [29/10] Năm gap theo ba chủ đề. Bỏ G1 legacy cũ: phần về oracle đã nằm ở G2 (WCA4Z, PL/I), G4 (IP-preserving), G5 (K-COBAgent, Hans et al.); 5 bài lại lẫn với track, nên đưa sang hướng nghiên cứu 3. Bỏ G6 cũ (bài không nêu hạn chế của oracle) vì là thiếu sót báo cáo, đưa vào Checklist (10 bài sau khi mã lại Google migrations).
\new{The five gaps concern three topics: oracles that do not execute the code (G1, G2); the origin and independence of tests (G3, G4); and the original program as a reference (G5).}

\begin{enumerate}[label=\textbf{G\arabic*.}]
  % [29/10] Từ G5 cũ, gộp thêm similarity (bản cũ để ở Implications). Không nhắc khoảng 3 đến 98 điểm đã có trong Hình 3; nêu số phải tự đếm: 13/18 bài lệch hơn 10 điểm, similarity đảo thứ hạng ở 4/5 bài (Bảng 9). Bỏ ví dụ ablation của DepTrans vì không phải kết quả chính. Thêm 9 bài không có oracle chạy code, 3 bài không có phép kiểm nào khác ngoài biên dịch hay độ giống (đã kiểm: VIM-PT chỉ BLEU-4; CodePlan build và Diff BLEU; Collaboration không nói "success rate" đo bằng gì). Pilot: 15/32 bản trượt test vẫn biên dịch được; CodeBLEU vừa nhận bản trượt vừa bác bản qua ở cả ba ngưỡng (0.5: 17/14; 0.6: 13/35; 0.7: 5/52).
  \item \textbf{\new{Compile rates and similarity scores diverge from behavioral results.}}
  \new{Compilation is the most common oracle, yet neither it nor similarity checks behavior (Table~\ref{tab:oracles}). Nine papers use no oracle that executes the code (Section~\ref{sec:rq1}), and 3 of them describe no check other than compilation or similarity~\cite{vimpt,codeplan,collaboration}. Where other papers report these checks next to a behavioral rate on the same outputs, they diverge from it (Section~\ref{sec:rq3}). The compile rate exceeds the behavioral rate by more than 10 percentage points in 13 of 18 papers, and similarity reverses at least one ranking of systems or settings in 4 of the 5 papers that compare it with a pass rate. On the pilot's single set of translations, compilation accepts nearly half of those that fail the tests; at each of its three thresholds, CodeBLEU both accepts some that fail and rejects some that pass (Section~\ref{sec:pilot-results}). Compilation still catches static errors and similarity measures closeness to a reference, but neither shows by itself that an output behaves like the original. A compile rate or a similarity score should therefore come with a behavioral rate on the same outputs.}
  % [29/10] Từ G4 cũ. Bỏ F1 của Moumoula et al. (clone detection, ngoài map, không đo độ đúng của migration); thay bằng mục 5.5 và pilot. Thêm 6 bài còn lại trong 9 bài không chạy code (Rule synthesis, MSA, Pig, PL/I, WCA4Z, DepUpgrade; đều có O-i hoặc O-h); LIB 9/12 có người chấm; Google: "quality is assured through the manual review process". 3/17 bài báo cáo độ đồng thuận. Pylibmig: κ trên 194 rồi 114 thay đổi, sửa hướng dẫn giữa hai đợt (0.83, 0.73), sau đó một người chấm phần còn lại; Fruntke & Krinke: hai kỹ sư chấm mù, đồng thuận theo từng tiêu chí. Hai bài O-h không có số, không so với oracle chạy code. Pilot: judge bác 67/164 lời giải tham chiếu, nhận 16/32 bản trượt. Khuyến nghị không quy định κ. Đã kiểm với bài gốc.
  \item \textbf{\new{Judgments by people or LLMs are rarely validated.}}
  \new{The other 6 papers without an oracle that executes the code use human or LLM judgment~\cite{rulesynth,msa,pig,pli,wca4z,depupgrade}. Such verdicts can serve as evidence of behavior, but only as far as their consistency or accuracy is known. Human evaluation is common in library migration, where 9 of the 12 papers include it (Section~\ref{sec:rq2}). Google's industrial migrations depend on it as well: LLM-generated changes go to the engineers who own the affected code, and quality ``is assured through the manual review process''~\cite{googlemig}. Yet only 3 of the 17 papers with human evaluation report agreement between raters (Section~\ref{sec:rq3}). Two of them show how, by revising the coding guideline between early batches until the reviewers agree substantially~\cite{pylibmig} or by reporting agreement for each criterion on a blind sample~\cite{fruntke}. Neither main-corpus paper with an LLM judge (O-h)~\cite{depupgrade,wca4z} gives a figure for the judge's accuracy or compares it with an oracle that executes the code, although WCA4Z relies on its judge alone. That accuracy can change with the programs being assessed (Section~\ref{sec:i3}). The pilot measures it: on answers known in advance, a small LLM judge rejects 41\% of the benchmark's own reference solutions; on the translations, it accepts half of those that fail the tests (Section~\ref{sec:pilot-results}). Human judgments should therefore come with agreement between raters, and LLM judges with a check on known answers from the domain where they will be used.}  
  % [29/10] Từ G2 cũ. Tiêu đề cũ ("limited demonstrated transfer") nói điều map không đo; tiêu đề mới là điều map đếm được (12/20 bài dịch chỉ dùng test benchmark; 6/7 bài LIB, VER chạy test của dự án). Thêm bằng chứng test có thể sót lỗi: 14/28 bài tự nêu test yếu (8 benchmark, 5 dự án, 1 cả hai; Bảng 10); pilot sót một số mutant, tới 16% trong 288 mutant biên dịch được; pylibmig: "only 27 migrations have at least one test that covers the code changes" (đã kiểm bài gốc). Thêm khuyến nghị.
  \item \textbf{\new{Test evidence for translation rests largely on benchmarks.}}
  \new{Tests execute the code, but in translation between languages they mostly come from benchmarks rather than from real projects. Twelve of the 20 translation papers with existing tests rely only on benchmark tests, whereas in library and version migration 6 of the 7 run the project's own suite (Section~\ref{sec:rq2}). A pass rate on benchmark problems need not carry over to code that involves persistent state, dependencies or interaction with the rest of an application. Since we did not extract test counts or coverage, the map does not measure how strong the tests are. The papers' own limitations and the pilot, however, show that benchmark tests and project suites can both miss faults. Of the 28 papers with existing tests, 14 name weak or few tests as their main limitation: 8 with benchmark tests, 5 with project suites and 1 with both (Section~\ref{sec:rq3}). In the pilot, the HumanEval-X tests miss some of the seeded faults, up to 16\% of those that compile (Section~\ref{sec:pilot-design}). Project suites may not even reach the migrated code: in Python library migration, only 27 of the 314 migrations have a test that covers the changed code~\cite{pylibmig}. Studies should therefore state where their tests come from and how much of the migrated code the tests cover.}
  % [29/10] Từ G3 cũ. Tiêu đề thành nhận định; câu 2 nêu hai kiểu phụ thuộc. 4/7 bài để hệ thống dịch cả code lẫn test (bài thứ năm trong bản cũ, IntentTester, chỉ migrate test nên test chính là output). Bỏ Syzygy (ngoài map), thay bằng ReCodeAgent: test của developer do người dịch hoặc kiểm, chỉ chạy sau khi agent xong (4.2.2). Ba bài same-source: IP-preserving, K-COBAgent, AlphaTrans. InterTrans, SmartC2Rust chọn hoặc sửa bằng chính test đánh giá; SmartC2Rust và FreshBrew (Case Study 1) có ví dụ lách test. RISE: lúc dịch và lúc chấm dùng cùng phép so kết quả, nhưng lúc chấm còn kiểm tay (2.1.3, 4.1.2). TRAVEL (MCTS chỉ để tạo dữ liệu train) và DepTrans thưởng bằng loại phép kiểm mà bài báo cáo. Cột cờ thiếu Google, K-COBAgent, nên viết "there may be more".
  \item \textbf{\new{Some tests are not independent of the system under evaluation.}}
  \new{Evidence of behavior should not depend on the system that it assesses. Tests can depend on the system in two ways: it produces them, or it uses them to shape its outputs. In 4 of the 7 papers with translated tests, the system under evaluation translates the code as well as its tests, so one translation error can affect both (Section~\ref{sec:rq2}). The map counts papers with this setup, not how often it leads to such errors. ReCodeAgent, one of these papers, also reports results on developer tests that were translated or checked by humans and that its agents never see, to guard against errors in the tests it translates itself~\cite{recodeagent}. Three of the 49 papers state as their main limitation that the oracle and the output come from the same source (Table~\ref{tab:limits}). IP-preserving, for instance, generates its tests and code from one specification~\cite{ipmodern}. The second kind arises when the system selects or repairs its outputs with the tests that later judge them. InterTrans stops its search once a candidate passes the benchmark tests on which it reports accuracy~\cite{intertrans}, and SmartC2Rust repairs its translations against the tests on which it is then evaluated~\cite{smartc2rust}. Their pass rates show that the outputs satisfy the tests used to select or repair them, not that they would pass independent ones. A model that sees the tests can even game them, as when SmartC2Rust's LLM simplifies code just to make the tests pass on a large program or an agent in FreshBrew excludes failing tests from the build~\cite{smartc2rust,freshbrew}. RISE likewise checks its translations while producing them by comparing the results of the source and translated queries, the same differential test that it reports. Its evaluation, however, adds a step that the technique does not apply to its outputs: the authors confirm by hand each translation whose results match~\cite{rise}. In training, TRAVEL and DepTrans reward their models with the kinds of check they report, a weaker form of this dependence, since they train on data distinct from their benchmarks~\cite{travel,dgir}. We noted these uses in search, repair and training as we came across them during extraction, so there may be more. Studies should therefore state whether the system under evaluation produced or used the tests, or was trained against checks of the same kind, and, if so, also report results on tests that it neither produced nor used.}
  % [29/10] Gap mới: tiêu chuẩn ở đoạn mở đòi so với bản gốc. 14/49 bài dùng differential testing; 4/28 bài có test sẵn cũng chạy bản gốc. K-COBAgent: 82/544 module có oracle dùng được, 462 còn lại gồm copybook, chương trình phụ thuộc CICS, không standalone, interactive (bài không tách nhóm nào thuộc 261 SKIPPED hay 201 PARTIAL). Hans et al.: đội kiểm thường thiếu database, hệ file nên mock, chạy trên mainframe. 6/14 bài O-f nêu hạn chế là phép so chưa đủ. "does not record": map không mã hóa lý do, còn IP-preserving có nêu (4.3). Bỏ WCA4Z ("impractical" là do bản dịch hiếm khi biên dịch được). Đã kiểm với bài gốc.
  \item \textbf{\new{Most studies do not run the original program.}}
  \new{The most direct evidence that migrated code behaves like the original is differential testing: running both on the same inputs and comparing their outputs. Only 14 of the 49 papers use it; of the 28 papers with existing tests, 4 also run the original (Section~\ref{sec:rq1}). Running it is not always possible. K-COBAgent runs the COBOL programs it migrates on generated inputs, but only 82 of its 544 modules yield a usable oracle; the other 462 include copybooks and programs that depend on CICS, are not standalone or are interactive~\cite{kcobagent}. Hans et al.\ note that validation teams often lack the databases and file systems COBOL programs need, so the authors mock these resources and run the programs on a mainframe~\cite{coboljava}. Where the original does run, the comparison covers only the inputs tried: 6 of the 14 papers name incomplete differential, fuzzing, symbolic or bounded checks as their main limitation (Table~\ref{tab:limits}). For the 35 papers that do not run the original, the map does not record whether it could not be run or the authors chose not to. Studies should therefore compare their outputs with those of the original where it can be run, alongside any tests, and otherwise state what prevented it, as K-COBAgent and IP-preserving do~\cite{kcobagent,ipmodern}.}
\end{enumerate}

% [29/10] Thay bốn hướng và đoạn MBTA của bản cũ: các hướng cũ là yêu cầu cho benchmark và cách báo cáo (trùng Implications, Checklist) và dẫn Syzygy, Moumoula et al. ngoài map; MBTA đo độ bền của translator, không phải oracle. Nay là ba câu hỏi nghiên cứu, mỗi câu gồm: vì sao còn mở (trỏ về G), đã có gì, cách nghiên cứu, kết quả dùng vào đâu (cả hai chiều).
\paragraph{\new{Directions for further study.}}
\new{Three questions for research on migration oracles follow from the gaps; for each, we note how it could be studied and what an answer would make possible.}
\begin{itemize}
  % [29/10] Ít nhất 18/49 bài phân loại lỗi của output (Lost in Translation, UniTrans, RustRepoTrans, TransLibEval, BatFix, SmartC2Rust, DepWareTrans, PROGnosticator, Pig, IntentTester, Rule synthesis, FreshBrew, RISE, GUIMigrator, Java-to-ArkTS, pylibmig, Translation Tag Team, Code hints). 5 bài mô tả lỗi lọt qua một phép kiểm nhưng bị phép khác bắt, mỗi bài nhiều nhất 3 loại lỗi và 3 oracle (Code hints: fuzzing so với review tay). Pilot: 25/32 bản trượt được ít nhất một oracle không chạy code nhận, chỉ 4/32 được cả ba nhận (ngưỡng 0.6). Đã đọc cả 49 bài.
  \item \textbf{\new{Which faults does LLM-based migration introduce, and which oracles catch them? (G1--G3)}}
  \new{Many oracles let faulty migrations through, as G1--G3 show for compilation, similarity, LLM judges and tests. What one oracle misses, another may catch. In the pilot, the three oracles that do not execute the code seldom all accept the same failing translation (Section~\ref{sec:pilot-results}). Kinds of fault are widely documented: at least 18 of the 49 papers classify the failures of migrated code, from translation between languages (Lost in Translation) to library and version migration (Pig, FreshBrew)~\cite{lostintranslation,pig,freshbrew}. Five papers also describe faults that slip past one check but not another~\cite{codehints,pylibmig,terramod,freshbrew,guimigrator}. In Code hints, for instance, some refactorings pass a fuzzing-based equivalence check yet fail manual review~\cite{codehints}. These accounts cover at most 3 kinds of fault and 3 oracles per paper, so no paper systematically measures what each oracle misses. The pilot already seeds generic faults into reference solutions to assess its tests and judge (Section~\ref{sec:pilot-design}). A study could instead seed faults of each kind found in these classifications into target programs known to be correct and run every oracle on them, recording which kinds each oracle lets through. Researchers would then know which oracles complement each other, and teams where their own oracles leave gaps.}
  % [29/10] RustRepoTrans Table III (tập con 300 bài CodeTransOcean, chương trình đứng riêng; Table I ghi Program Level): Pass@1 giảm 16.2 đến 30.8 điểm; V3 kém Claude-3.5 8.0 điểm (66.3 so với 74.3) nhưng hơn 6.6 (50.1 so với 43.5). Không viết "model đứng đầu đổi" vì Claude hơn R1 chỉ 0.6 điểm. ReCodeAgent 4.2.2: test tool tự dịch qua 1,456/1,484 (không gồm Crust), test ẩn qua 1,822/2,107; tác giả quy một phần chênh lệch cho test ẩn dịch sai (csv, 57 lần trượt). Trong 49 bài, chỉ RustRepoTrans báo thứ hạng đổi (Lost in Translation: GPT-4 đứng đầu ở cả hai; DepUpgrade giữ thứ tự), chỉ ReCodeAgent báo pass rate trên cả hai loại test. RustRepoTrans cũng chọn dự án đã được viết lại sang Rust và dùng test của bản Rust, nên thiết kế làm được.  
  \item \textbf{\new{How far do published results carry over to real projects? (G3, G4)}}
  \new{Most translation papers with existing tests take them only from benchmarks (G3), and some use tests that the system under evaluation produced or used (G4). Results gathered this way may not carry over to real projects, either because real code differs from benchmark problems or because a system can score higher on tests that it produced or used than on independent ones. RustRepoTrans shows how much the first can matter: the pass rates of 4 LLMs fall by 16 to 31 percentage points between a benchmark of stand-alone programs and a set of functions from real projects that depend on other code in their repository~\cite{rustrepotrans}. DeepSeek-V3 scores 8 points below Claude-3.5 on the benchmark but 7 points above it on the real projects, so the ranking of the models changes. As for the second, ReCodeAgent's translations pass 98.1\% of the tests that the tool translated itself but only 86.5\% of hidden developer tests~\cite{recodeagent}. The two sets differ, however, and the authors trace some of the hidden failures to errors in how those tests were translated, so the gap need not mean that the tool's tests favor it. Among the 49 papers, the change in ranking appears only in RustRepoTrans and a comparison of pass rates on the two kinds of test only in ReCodeAgent, so neither has been examined for other systems or beyond translation. Researchers could collect real projects that their developers have already migrated, for instance to another language or to a new library, version or framework, and hold back the tests of the migrated versions for scoring, or part of them for systems using tests while migrating. Running several LLMs and published migration tools on these projects and on existing benchmarks would reveal how often the systems rank differently on the benchmarks than on the projects. Because the held-back tests pass on the developers' own migration, they would also show whether a system's pass rate on tests that it produced or used can be trusted. Teams planning a migration would then know when published results can guide their choice of system, and researchers would see which benchmarks rank systems in the same order as real projects.}
  % [29/10] IP-preserving 4.3: cần "a production-like TOS instance, recorded traffic, and a data-protection regime permitting replay"; test sinh từ cùng đặc tả, chỉ kiểm trên dữ liệu tổng hợp. K-COBAgent: GnuCOBOL, bỏ qua chương trình không chạy riêng được; Table 1: 82 PASS hoặc FAIL, 201 PARTIAL, 261 SKIPPED. Hans et al.: symbolic execution theo paragraph, mock database, file, CICS (cả phía Java), báo branch coverage; không bài nào so với chạy ở môi trường đầy đủ. TerraMod chạy trên LocalStack thay AWS. Không viết GnuCOBOL khác trình biên dịch thật (không rõ chương trình X-COBOL viết cho trình biên dịch nào).  
  \item \textbf{\new{How much is lost when the original program runs outside its full environment? (G5)}}
  \new{Running the original and the migrated code on the same inputs is the most direct way to check whether behavior is preserved, but the original may need an environment that a study lacks (G5). This matters especially in legacy migration, where none of the 5 papers uses existing tests (Section~\ref{sec:rq2}), so a study that does not run the original falls back on references such as a specification, an LLM judge or people (G2, G4). IP-preserving, which tests its code against a specification, names what running the original would require: a production-like instance of the legacy system, recorded traffic and data-protection rules that permit replaying it~\cite{ipmodern}. Only 2 of the 5 legacy papers run the original, and both do so without its full environment. K-COBAgent skips the programs that cannot run on their own, and most of its 544 modules are never compared with the original (G5)~\cite{kcobagent}. Hans et al.\ run each COBOL paragraph on inputs from symbolic execution and replace its database, file and CICS calls with mocks~\cite{coboljava}. Both report how much of the code their comparisons reach, as the number of modules compared or as branch coverage, but checking whether these comparisons catch the same faults as running the original on recorded traffic would need the full environment that they work around. Where the full environment is still available, for instance at a company that operates the legacy system, researchers could compare the same migrated programs with the original both there on recorded traffic and outside it on generated inputs with mocked resources. Recording the faults the second comparison misses, the correct programs it rejects and the programs it cannot handle would show how much is lost and where. Tests that run the migrated code with mocks or emulators, as TerraMod does on an emulated AWS cloud~\cite{terramod}, could likewise be compared with runs in the full environment. Teams migrating systems that are hard to run would then know whether checking outside the full environment is enough or setting it up is worth the effort, and researchers would see which kinds of program cannot yet be handled with mocks, emulators or generated inputs.}
\end{itemize}

\subsection{Implications}
\label{sec:implications}

% [29/10] Bỏ hai câu khuyến nghị trùng G1, G3, G4 (đã có ở cuối mỗi gap) và câu many-shot (ngoài map, nói về độ dài prompt; bài gốc không đỡ cho ý của câu). Đoạn chỉ trỏ sang Checklist (tác giả theo, reviewer kiểm) và Directions. Checklist phải có mục cho từng khuyến nghị của G1–G5.
\paragraph{For researchers.}
\new{Each of the five gaps (G1--G5) ends with a recommendation for studies that evaluate migration; the checklist in Section~\ref{sec:checklist} gathers these recommendations into a list that authors can follow and reviewers can check. Questions that better reporting cannot answer, such as which oracles complement each other, are taken up in the directions for further study (Section~\ref{sec:rq5}).}

% [29/10] Bỏ số similarity và ví dụ TerraMod (G1 đã nói), "reproducible environments" (hướng 3 đã bàn môi trường thay thế) và câu nhóm thưa (legacy ở hướng 3; map không cho thấy nhóm kia thiếu benchmark). Thêm: giữ đủ ca khi chuyển test (pilot: HumanEval/131 thiếu ca 2468 mà test Python có); mutation check; mutant bị test bác làm đáp án sai để kiểm judge (pilot); giữ một phần test ẩn (InterTrans dừng tìm kiếm ở ứng viên đầu tiên qua test benchmark; agent trong FreshBrew loại test khỏi build ở Case Study 1, cho test tự bỏ qua ở Case Study 3). FreshBrew: "many reward hacking instances also occur below this threshold" (5 điểm coverage). Đã kiểm với bài gốc.
% [29/10] Bảng checklist (mục 7.3) đặt ở đây để LaTeX đưa nó lên đầu trang ngay trước mục 7.3; đặt trong 7.3 thì bảng trôi sang mục 8.
\begin{table}[t]
\centering
\caption{\addedtag Oracle Provenance Checklist. Gap: the gap in Section~\ref{sec:rq5} that the item addresses. In the map: how many main-corpus papers report the item, where our coding records it.}
\label{tab:checklist}
\small
\begin{tabular}{@{}l P{0.48\linewidth} l P{0.31\linewidth}@{}}
\toprule
& Report & Gap & In the map \\
\midrule
1 & Each oracle (Table~\ref{tab:oracles}) and the behavior it observes, such as return values, output, state, termination or the rendered screen; the population and denominator of every rate & --- & Not coded; 2 of the 18 papers in Figure~\ref{fig:compile-gap} give the behavioral rate over compiling outputs only \\
\addlinespace
2 & A behavioral rate on the same outputs next to every compile rate or similarity score & G1 & Compilation: 18 of 35 papers; similarity: 5 of 13 \\
\addlinespace
3 & For human judges, the raters, their independence, the rubric, the sample and the agreement between raters; for LLM judges, the model version, the prompt and the accuracy on known correct and incorrect answers from the same domain & G2 & Rater agreement: 3 of 17; judge accuracy: 0 of 2 \\
\addlinespace
4 & Who wrote the tests and where their expected outputs come from (benchmark, project developers, the authors, the system, an LLM or a specification); the number of tests and their coverage of the migrated code & G3 & Origin not described for some oracle in 3 of the 40 papers with test-based oracles; counts and coverage not coded \\
\addlinespace
5 & Whether the system under evaluation produced the tests, selected or repaired its outputs with them, or was trained against the same kind of check; if so, results on tests that it neither produced nor used & G4 & Not coded systematically; held-out results in 1 paper~\cite{recodeagent} \\
\addlinespace
6 & A comparison with the original program on the same inputs, with the source of the inputs, or otherwise what prevented running the original & G5 & Original run: 14 of 49 \\
\addlinespace
7 & The oracle's known failure modes in the evaluated setting, such as unexercised paths, timeouts, flaky tests or ways a system can game it & --- & Stated: 39 of 49 \\
\bottomrule
\end{tabular}
\end{table}

\paragraph{For benchmark builders.}
\new{Many studies take their tests from benchmarks (G3), so builders can make these tests stronger and harder to game. When porting tests to another language, they could keep every case: in the pilot, one seeded fault survives because the Java tests of its HumanEval-X problem lack a case that the Python tests check (Section~\ref{sec:pilot-design}). A mutation check like the pilot's would flag such weak tests before release. The mutants that the tests reject and the reference solutions would also provide known incorrect and correct answers on which to check LLM judges (G2). Builders could keep part of the tests hidden from the systems they evaluate (G4). Visible tests can be gamed either by a system that selects or repairs its outputs against them, as InterTrans does in its search~\cite{intertrans}, or by an agent that excludes failing tests from the build or makes a test skip itself, as some agents do in FreshBrew~\cite{freshbrew}. FreshBrew's own coverage check catches only part of this second kind.}

% [29/10] Bỏ các câu chung ("Select checks…", "Acceptance criteria…"), câu mock và workload tổng hợp (hướng 3 đã nói) và vế "evaluations should document" (việc của tác giả, G5 đã khuyến nghị); câu Google chỉ thuật lại việc họ làm nên chuyển thành việc team làm được. Đoạn theo thứ tự việc của một team. Chọn tool: kiểm oracle và nguồn test của kết quả công bố (G1–G3; "need other evidence" theo G1 "neither shows by itself" và G2 "only as far as their consistency or accuracy is known"), rồi thử danh sách đã lọc trên một phần nhỏ của hệ thống cần migrate, chấm bằng test tool không thấy (G3, G4). Lúc migrate: kiểm thay đổi agent làm với test và build (FreshBrew Case Study 1 sửa pom.xml, Case Study 3 sửa test); đo đồng thuận giữa người review khi dựa vào review như Google (G2). Kiểm với bản gốc (G5): môi trường, traffic đã ghi, quyền phát lại theo IP-preserving 4.3; "may need" vì input sinh tự động cũng dùng được.
\paragraph{For practitioners migrating real systems.}
\new{A team should check which oracle and which tests a published result rests on before using it to choose a tool, since oracles that do not execute the code need other evidence to show whether outputs behave like the original (G1, G2) and pass rates on benchmark problems may differ on a real project (G3). With a shortlist in hand, the team can find a suitable tool by running each one on a small part of the system to be migrated and scoring the results with tests the tools never see (G3, G4). An agent's changes to the tests or the build during the migration should be inspected, since agents in FreshBrew used such changes to sidestep failing tests~\cite{freshbrew}. Where the team relies on code review, as Google's migrations do~\cite{googlemig}, the agreement between reviewers who judge the same sample of outputs can show how reliable the review is (G2). Comparing the migrated code with the original gives the most direct evidence (G5) but may need a production-like environment, recorded traffic and permission to replay it; these are easier to secure before the old system is retired.}

\subsection{A checklist for oracle provenance}
\label{sec:checklist}

% [29/10] Viết lại checklist: bảy mục bám theo G1--G5 thay cho ba tầng cũ; ký hiệu O-x như cả bài (bỏ $O_b$). Mục 6 (so với chương
% trình gốc) là khuyến nghị của G5, bản cũ không có; mục 7 thay G6 cũ (bài không nêu hạn chế, 10 bài sau khi mã lại Google migrations).
% Cột "In the map" lấy từ mục 5: 18/35 (Hình 3), 5/13 (Bảng 9), 3/17 và 0/2 (5.4), 3/40 bài có oracle dựa trên test để trống nguồn
% gốc (TransAgent e:U; K-COBAgent, ReCodeAgent f:U; Bảng 7), 14/49 (O-f), 39/49 (Bảng 10); "not coded" khi map không trích trường đó.
\begin{added}
A reported success rate can be interpreted only if readers know which oracle produced it, where the oracle's ground truth comes from and how far the oracle has been checked. Table~\ref{tab:checklist} turns the recommendations that end G1 to G5 into seven items that a study of LLM-based migration can report and a reviewer can check. Items~2 to~6 follow the gaps in order. Item~1 asks for the definition behind every rate, since equivalence always holds with respect to some oracle and some notion of observable behavior (Section~\ref{sec:background}). Item~7 asks for the limits of the oracle, which 10 of the 49 papers do not state (Table~\ref{tab:limits}).

The last column gives a baseline from the map where our coding records the item. It counts what papers report, not what they did, so it shows how far current reporting is from the checklist rather than how good individual studies are. The checklist sets no thresholds, such as a minimum agreement or coverage, because suitable values depend on the system and on the risk of the migration. Nor does it require a particular oracle: a study that cannot run the original program meets item~6 by stating what prevented it.
\end{added}
